% \numberwithin{equation}{section}

\chapter{Ones MHD} \label{sec:onesMHD}
En la secció anterior s'han presentat totes les equacions rellevants amb la descripció de cada paràmetre que apareix en elles.
Les equacions que s'han de linealitzar de l'apartat anterior per a la discussió de les ones són les següents:
llei de Gauss, Eq.~(\ref{eq.2.2}); inducció, Eq.~(\ref{eq.2.7}); continuïtat de massa, Eq.~(\ref{eq.2.8}); equacions del moviment, Eq.~(\ref{eq.2.10}); i calor, Eq.~(\ref{eq.2.14}) \footnote{La temperatura, $T$, és una variable que es pot calcular a partir de la densitat, $\rho$, i de la pressió, $p$, llavors no intervé directament en les equacions.}.
Aquest procés és semblant al duit a terme en els articles de \textcite{oscillation_modes_II} i d'\textcite{oscillations_quiescent_solar_prominence},
deixant de banda la deducció més general del llibre de \textcite{Priest_2014}.


\section{Linealització} \label{subsec:linealització}
Per linealitzar les equacions s'ha de descomposar cada terme en una part estacionària o quiescent i en una part pertorbativa, suposadament molt menor.
D'aquesta manera les variables són del següent estil:
\begin{equation}
	f(\bm{r},t) = f_0 + f_1(\bm{r},t) \quad , \quad |f_0| \gg |f_1(\bm{r},t)|. \tag{3.1} \label{3.1}
\end{equation}

Les que es fan servir i es substitueixen en les equacions mencionades són la densitat, la pressió, la velocitat i el camp magnètic:
\begin{gather}
	\rho = \rho_0 + \rho_1, \tag{3.2} \label{3.2} \\
	p = p_0 + p_1, \tag{3.3} \label{3.3} \\
	\bm{v} = \bm{v}_1, \tag{3.4} \label{3.4} \\
	\bm{B} = \bm{B}_0 + \bm{B}_1. \tag{3.5} \label{3.5}
\end{gather}

La condició d'equilibri estàtic (\hyperref[punt:I]{I}) imposa una velocitat quiescent nul·la ($\bm{v}_0 = \bm{0}$).
A l'hora de linealitzar només es tenen en compte els termes d'ordre zero (variables quiescents i els seus productes) i els termes de primer ordre, les pertorbacions lineals.
Llavors s'ignoren els ordres superiors donats pels productes de les parts pertorbatives de dues variables o el quadrat d'una mateixa.

Les cinc equacions citades en l'inici es linealitzen de la manera indicada, llavors els resultats són els següents:
\begin{gather}
	\frac{\partial \rho_1}{\partial t} + (\bm{v}_1 \cdot \bm{\nabla}) \rho_0 + \rho_0 (\bm{\nabla} \cdot \bm{v}_1) = 0, \tag{3.6} \label{eq.3.6} \\
	\rho_0 \frac{\partial \bm{v}_1}{\partial t} = - \bm{\nabla} p_1 + (\bm{\nabla} \bm{\times} \bm{B}_1) \bm{\times} \bm{B}_0 / \mu, \tag{3.7} \label{eq.3.7} \\
	\frac{\partial p_1}{\partial t} + (\bm{v}_1 \cdot \bm{\nabla}) p_0 - c_s^2 \left( \frac{\partial \rho_1}{\partial t} + (\bm{v}_1 \cdot \bm{\nabla}) \rho_0 \right) = 0, \tag{3.8} \label{eq.3.8}
\end{gather}
\begin{gather}
	\frac{\partial \bm{B}_1}{\partial t} = \bm{\nabla} \bm{\times} (\bm{v}_1 \bm{\times} \bm{B}_0), \tag{3.9} \label{eq.3.9} \\
	\divBp = 0, \tag{3.10} \label{eq.3.10}
\end{gather}
on la velocitat del so, $c_s$, ve donada per
\begin{equation}
	c_s^2 = \frac{\gamma p_0}{\rho_0} = \frac{\gamma k_B T_0}{m}, \tag{3.11} \label{eq.3.11}
\end{equation}
i es dedueix a partir de la substitució de l'Eq.~(\ref{eq.3.6}) dins l'Eq.~(\ref{eq.3.8}).

Per linealitzar l'equació de calor, Eq.~(\ref{eq.2.14}), i arribar a l'Eq.~(\ref{eq.3.8}) s'ha hagut de fer una aproximació en sèrie de Taylor de l'estil
\begin{equation}
	\frac{1}{\rho^{\gamma}} = \frac{1}{\rho_0^{\gamma}} \left( 1 + \frac{\rho_1}{\rho_0} \right)^{-\gamma} \approx \frac{1}{\rho_0^{\gamma}} \left( 1 - \gamma \frac{\rho_1}{\rho_0} \right), \tag{3.12} \label{eq.3.12}
\end{equation}
apart de fer servir la definició de la velocitat del so de l'Eq.~(\ref{eq.3.11}).

Per resoldre les equacions linealitzades es proposen ones planes ja que el medi és uniforme i es troba en equilibri quiescent.
Les ones planes impliquen una dependència de les variables d'espai i temps com a exponencials complexes
\begin{equation}
	\bm{f}_1(\bm{r},t) = \bm{f} e^{i(\bm{k} \cdot \bm{r} - \omega t)}. \tag{3.13} \label{eq.3.13}
\end{equation}

La propagació de les pròpies ones té lloc la direcció $\hat{z}$.
Per la simetria del sistema (Fig.~\ref{fig:sistema}), tant les variables pertorbades escalars de densitat, $\rho$, i pressió, $p$, com les vectorials de velocitat, $\bm{v}$, i
camp magnètic, $\bm{B}$, representades per $f$ a l'Eq.~(\ref{eq.3.13}), només depenen de la posició, $x$, en que s'avaluen, llavors
\begin{gather}
	\rho_1(\bm{r}, t) \equiv \rho(x) e^{i(\omega t - k_z z)}, \tag{3.14} \label{3.14} \\
	p_1(\bm{r}, t) \equiv p(x) e^{i(\omega t - k_z z)}, \tag{3.15} \label{3.15} \\
	\bm{v}_1(\bm{r}, t) \equiv \big(v_{x}(x), v_{y}(x), v_{z}(x)\big) e^{i(\omega t - k_z z)}, \tag{3.16} \label{3.16} \\
	\bm{B}_1(\bm{r}, t) \equiv \big(B_{x}(x), B_{y}(x), B_{z}(x)\big) e^{i(\omega t - k_z z)}. \tag{3.17} \label{3.17}
\end{gather}
Es fa evident que hi ha vuit funcions desconegudes: $\rho(x)$, $p(x)$, les sis components dels vectors $\bm{B}(x)$ i $\bm{v}(x)$, juntament amb les incògnites $\omega$ i $k_z$.
Aquestes són les que s'empren per calcular els autovalors i les autofuncions dels modes normals, apart de la relació de dispersió.

Les derivades actuen com operadors ja que es tracta d'ones planes, del tipus de l'Eq.~(\ref{eq.3.13}).
Llavors aquestes es transformen de la següent manera:
\begin{equation}
	\frac{\partial}{\partial t} \implies i \omega, \tag{3.18} \label{eq.3.18}
\end{equation}
\begin{equation}
	\bm{\nabla} \implies \hat{x} \frac{\partial}{\partial x} - i k_z \hat{z}. \tag{3.19} \label{eq.3.19}
\end{equation}

Es substitueixen els operadors dins les equacions linealitzades i s'arriba a unes noves equacions més senzilles.
L'objectiu final de les equacions és descriure el sistema segons les autofuncions, $\bm{v}$, els autovalors, $\omega$, i les constants que el caracteritzen.

A partir de l'Eq.~(\ref{eq.3.6}) s'obté la següent expressió:
\begin{equation}
	i \omega \rho + \rho_0 \left( \frac{\partial \tilde{v}_{x}}{\partial x} - i k_z v_{z} \right) = 0, \tag{3.20} \label{eq.3.20}
\end{equation}
i si es defineix ${\tilde{v}_{x}}(x) \equiv i {v}_{x}(x)$ i es multiplica per $-i$ s'obté una equació sense la constant~imaginària,~$i$, explícitament
\begin{equation}
	\omega \rho - \rho_0 \left( \frac{\partial \tilde{v}_{x}}{\partial x} + k_z v_{z} \right) = 0. \tag{3.21} \label{eq.3.21}
\end{equation}
Aquesta es fa servir més envant, ja que es pretén eliminar la variable de densitat.

De la condició de camp magnètic amb divergència nul·la, Eq.~(\ref{eq.3.10}), s'aconsegueix la següent relació:
\begin{equation}
	\frac{\partial B_{x}}{\partial x} = k_z \tilde{B}_{z}. \tag{3.22} \label{eq.3.22}
\end{equation}

De l'equació d'inducció linealitzada, Eq.~(\ref{eq.3.9}), s'obtenen tres relacions ja que aquesta és vectorial.
De nou es defineixen les variables amb tilde ${\tilde{B}_{y}}(x) \equiv i {B}_{y}(x)$, ${\tilde{B}_{z}}(x) \equiv i {B}_{z}(x)$ per absoribir la constant imaginària, llavors
\begin{gather}
	\omega B_{x} = B_0 k_z v_{z}, \tag{3.23} \label{eq.3.23} \\
	\omega \tilde{B_{y}} = B_0 \frac{\partial v_{y}}{\partial x}, \tag{3.24} \label{eq.3.24} \\
	\omega \tilde{B_{z}} = B_0 \frac{\partial v_{z}}{\partial x}. \tag{3.25} \label{eq.3.25}
\end{gather}

Aleshores existeixen tres equacions pel camp magnètic estretament connectades amb els autovalors i les autofuncions.

De les equacions del moviment del plasma, Eq.~(\ref{eq.3.7}), també resulten tres noves relacions:
\begin{gather}
	\omega \rho_0 \tilde{v}_{x} = - c_s^2 \frac{\partial \rho}{\partial x}, \tag{3.26} \label{eq.3.26} \\
	\omega \rho_0 v_{y} = - \frac{B_0}{\mu} \frac{\partial \tilde{B_{y}}}{\partial x}, \tag{3.27} \label{eq.3.27} \\
	\omega \rho_0 v_{z} = c_s^2 k_z \rho + \frac{B_0}{\mu} \left( k_z B_{x} - \frac{\partial \tilde{B_{z}}}{\partial x} \right). \tag{3.28} \label{eq.3.28}
\end{gather}

Aquestes suposen més condicions que s'han d'acomplir juntament amb les tres anteriors.

S'han fet servir exponents complexos (ones planes) per treballar les equacions,
però s'ha de tenir en compte que les variables són reals i no imaginàries.
D'aquesta manera s'ha de prendre la part real de les ones planes: $\Re(f_1) = f (x) \cos(\omega t - k_z z)$.
Aquelles funcions representades amb tilde presenten un desfasament de $\pi /2$ respecte a les altres
\begin{align*}
	f = - i \tilde{f} \implies f_1 = - i \tilde{f_1} e^{i(\omega t - k_z z)} = \tilde{f_1} e^{i(\omega t - k_z z - \pi /2)}, \\
	\implies \Re(f_1) = \tilde{f_1} \cos(\omega t - k_z z - \pi / 2) \quad \text{on} \quad \tilde{f_1} \in \Re.
\end{align*}

Existeixen dos grups de variables desacoblades en aquest gran conjunt d'equacions.
En el primer bloc, en les Eqs.~(\ref{eq.3.24}) i (\ref{eq.3.27}) només apareixen les variables $v_{y}$ i $\tilde{B_{y}}$.
En canvi, en el segon bloc d'Eqs.~(\ref{eq.3.21}), (\ref{eq.3.23}), (\ref{eq.3.25}), (\ref{eq.3.26}) i (\ref{eq.3.28})
les variables que apareixen són les restants: $\rho$, $B_{x}$, $\tilde{B_{z}}$, $\tilde{v}_{x}$ i $v_{z}$.
Aquests dos blocs corresponen a les solucions dels modes d'Alfvén i als modes ràpids i lents, respectivament.



\section{Solucions: modes normals} \label{subsec:modes_normals}
Com s'ha comentat prèviament les variables principals són les de velocitat, $\bm{v}$, i camp magnètic, $\bm{B}$,
llavors s'expressen les equacions en funció aquestes.

Per obtenir els modes normals és necessari resoldre les equacions d'autovalors proposades prèviament.
Els dos blocs d'equacions es poden simplificar paral·lelament i resulten en els problemes d'Alfvén i de modes ràpids i lents respectivament.

Per resoldre els problemes i obtenir solucions del tipus modes normals s'han d'imposar les condicions de contorn de velocitat nul·la als extrems
\begin{equation}
	\bm{v} (x = \pm L) = \bm{0}. \tag{3.29} \label{eq.3.29}
\end{equation}
així com s'ha explicat en la Sec.~(\ref{sec:intro}).


\subsection{Modes d'Alfvén}
En aquest cas s'ha d'eliminar el camp magnètic de les Eqs.~(\ref{eq.3.24})~i~(\ref{eq.3.27}), el que resulta en una única equació d'autovalors:
\begin{equation}
	\omega^2  v_{y} + v_A^2 \frac{d^2v_{y}}{dx^2} = 0, \tag{3.30} \label{eq.3.30}
\end{equation}
on $v_A^2 \equiv \frac{B_0^2}{\mu \rho_0}$ és la definició de la velocitat d'Alfvén al quadrat.
Tant l'equilibri, com les equacions de contorn, Eq.~(\ref{eq.3.29}), com l'Eq.~(\ref{eq.3.30}) són simètrics respecte de $x=0$,
motiu pel qual les autofuncions dels modes d'Alfvén, $v_y(x)$, són simètriques o antisimètriques respecte de $x=0$.
Tot i que el problema d'Alfvén d'aquest cas té sol·lució analítica,
es resol amb Dedalus per conèixer el funcionament d'aquesta llibreria.


\subsection{Modes ràpids i lents} \label{subsubsec:ones_modes_r_i_l}
El segon bloc d'equacions és més complexe degut al gran nombre de variables i equacions.
Tot i així, la idea és la mateixa: eliminar les variables de densitat i camp magnètic per obtenir relacions d'autofuncions, $\tilde{v}_{x}$ i $v_{z}$, i d'autovalors, $\omega$.
En aquest cas hi ha dues equacions ja que hi ha dues velocitats en les variables. De nou, és una qüestió d'àlgebra el combinar les
Eqs.~(\ref{eq.3.21}), (\ref{eq.3.23}), (\ref{eq.3.25}), (\ref{eq.3.26}) i (\ref{eq.3.28}) per obtenir els següents resultats:
\begin{gather}
	\omega^2 \tilde{v}_{x} + c_s^2 \frac{d^2 \tilde{v}_{x}}{dx^2} + c_s^2 k_z \frac{d v_{z}}{dx} = 0, \tag{3.31} \label{eq.3.31} \\
	v_A^2 \frac{d^2 v_{z}}{dx^2} + k_z c_s^2 \frac{d \tilde{v}_{x}}{dx} + \left(-c_s^2 k_z^2 - v_A^2 k_z^2 + \omega^2 \right) v_{z} = 0. \tag{3.32} \label{eq.3.32}
\end{gather}
on $k_z$ és la constant d'acoblament. En el cas de que sigui nul·la, $k_z = 0$, es recupera el problema d'Alfvén per partida doble, per $\tilde{v}_{x}$ i per $v_z$.
En aquest cas, l'Eq.~(\ref{eq.3.31}) descriu els modes lents, amb velocitat pertorbada purament longitudinal, i l'Eq.~(\ref{eq.3.32}) descriu els modes ràpids, amb velocitat pertorbada purament vertical.
És precisament en aquesta situació que es calcula el comportament de cada mode ràpid i lent amb l'ajuda de Mathematica. Veure Apèndix~\ref{ap.B}.

Tal i com succeeix amb els modes d'Alfvén,
els modes ràpids i lents també tenen autofuncions $v_x(x)$ i $v_z(x)$ que són simètriques o antisimètriques respecte de $x=0$.
A més, l'estructura de les Eqs.~(\ref{eq.3.31})~i~(\ref{eq.3.32}) implica que si $\tilde{v}_{x}$ és parell respecte de $x = 0$, llavors $v_z$ és senar, i viceversa.

%---------------------------------------------------------------------------------------------------------------------------------------------------
% A partir d'observacions podem inferir es valors dels paràmetres del sistema - idealment - sismologia de protuberància

% CREC QUE AIXÒ M'HO VA DIR EN MONXO A COLQUE TUTORIA... NO SÉ SI HA D'ANAR AQUÍ O A UNA ALTRA BANDA
%---------------------------------------------------------------------------------------------------------------------------------------------------



\subsection{Relació de dispersió}
Per estudiar l'evolució dels autovalors de cada mode es representa el diagrama de dispersió.
Aquest ve donat per la relació entre la freqüència, $\omega$, que en aquest cas és l'autovalor, i el seu nombre d'ona, $k_z$.
El lligam d'aquests termes resulta en la velocitat de fase, $v_{ph}$, a la que es propaga l'ona
\begin{equation}
	{v}_{ph} = \frac{\omega}{k_z}. \tag{3.33} \label{eq.3.33}
\end{equation}
%La rellevància de la velocitat es
%
%
%Al resoldre les Eqs.~(\ref{eq.3.31}) i (\ref{eq.3.32}) s'obtenen els modes normals,
%un grup de solucions que s'han de treballar per arribar als modes ràpids i lents.
%El medi és dispersiu ja que hi ha un mode ràpid i té diferent veloctiat de fase
%Per aquest motiu al, mesclar modes es 

En el cas del problema d'Alfvén, Eq.~(\ref{eq.3.30}), no existeix cap relació de dispersió pel simple fet de que no apareix el nombre d'ona en la direcció $\hat{z}$, $k_z$.
Després de totes les consideracions preses sobre el sistema de la Fig.~\ref{fig:protuberancia},
els modes ràpids i lents només depenen del nombre d'ona, $k_z$, i de les velocitats de propagació del medi, $c_s$ i $v_A$,
com indiquen les Eqs.~(\ref{eq.3.31}) i (\ref{eq.3.32}).

En aquest treball es suposa que tant la protuberància com la corona són medis homogenis,
llavors les seves velocitats del so i d'Alfvén són constants a cada medi:
\begin{equation}
c_s(x) = 
\begin{cases}
c_{s p} & \text{per } |x| \leq x_p, \\
c_{s c} & \text{per } x_p < |x| \leq x_c,
\end{cases} \tag{3.34} \label{eq.3.34}
\end{equation}
\begin{equation}
v_A(x) = 
\begin{cases}
v_{A p} & \text{per } |x| \leq x_p, \\
v_{A c} & \text{per } x_p < |x| \leq x_c,
\end{cases} \tag{3.35} \label{eq.3.35}
\end{equation}
on $2x_p$ és l'amplada de la protuberància i $2x_c$ és l'amplada total del sistema, com es pot veure en la Fig.~\ref{fig:sistema}.

%---------------------------------------------------------------------------------------------------------------------------------------------------
% Crec que tot això no fa falta comentar molt en la memòria però ho hauria de tenir clar per la presentació.
%
% Modes Alfvén són intermitjos entre ràpids i lents... segons Priest, pàg. 155
% 
% Les ones d'Alfvén són del tipus Fig. 4.1.a del Priest, no?
% han de correspondre a l'Eq. 4.30? no sé si és un subcas o que
% 
% REALMENT NI UNES NI ALTRES JA QUE NO IGNORAM EL TERME DE CS^2
% 
% Repàs seccions 4.6, ones magneto-acústiques i 4.8 resum ones
%
%---------------------------------------------------------------------------------------------------------------------------------------------------
